\sscp{Words that straddle the Blust line}{15-04-02}
The discussion of **mantəR ‘cuscus’ in the previous section
and in \srf{14-04}, raises methodological challenges.
Its presence in around a dozen languages in two branches of \tsc{\ili{Muna}-Buton}
languages in Sulawesi as \it{mante} is too systematic to be dismissed as borrowing.
The medial *nt cluster is retained
and the final consonant is lost as expected.
And the same form or close variations distributed in many subgroups
across Wallacea also cannot be ignored.
Methodologically, if a word is found (preferably systematically)
in two or more lower-level subgroups,
then the common form is normally assigned to their shared parent language.
Based on our current understanding, that would be PMP.
But because marsupials are not found west or north of Sulawesi,
assigning **mantəR ‘cuscus’ to PMP would be unreasonable.\footnote{
		The only way **mantəR could be assigned to PMP would be if
		it has undergone semantic shift. However, even this
		would not explain the absence of cognates in other regions.}
So either the subgrouping assumptions
and boundaries need to be adjusted,
or other explanations need to be part of the discussion,
such as spread by diffusion or contact.
The distribution of **mantəR ‘cuscus’ is not bounded by the Blust line.
Nor is it found in all nodes or subgroups \ili{east} of the Blust line.
Both of these issues undermine its usefulness as evidence in support of CEMP.
And there is currently no Austronesian proto-language to which
it can be assigned that is based on exclusively shared innovations
in the historical phonology that accounts for the distribution
of its reflexes straddling the Blust line.

In \srf{13-01-05} we discussed PCEMP *keRa(nŋ) ‘hawksbill turtle’.
We noted that it links consistently to related forms in the Philippines,
western Indonesia and Madagascar meaning ‘shell, turtle shell’.
We observed that glosses of ‘shell,
turtle shell’ persist into Oceania; that glosses narrowed in
meaning to ‘hawksbill turtle’ are found in very few languages
in Wallacea; and that broader glosses relating to generic ‘turtle’
and other kinds of turtles are found throughout Wallacea
and persist into Oceania.
This term is more appropriately reassigned as PMP *kəRa(ŋ) ‘shell, turtle shell’.